\subsection{RNN-Like Models and State Space Models (SSM)}
\label{subsec:ssm}

Recurrent and state-space formulations offer a fundamentally different approach to sequence modeling: instead of computing pairwise interactions among all tokens (as in attention), they maintain a fixed-size hidden state that is updated incrementally at each time step, achieving $O(N)$ training complexity and $O(1)$ per-step inference cost. The core trade-off is between the \emph{expressivity} of the state update rule and the \emph{capacity} of the fixed-size state: a richer update rule can extract more information from each input token, but the state can store at most $O(d_s^2)$ bits regardless of the sequence length, where $d_s$ is the state dimension. This section traces the evolution of SSM and RNN-like token mixers from the foundational S4 model (2021) through the Mamba family of selective SSMs (2024--2025), modern RNN variants (RWKV, RetNet, xLSTM), and long-convolution architectures (Hyena, MEGA), culminating in the Structured State Space Duality (SSD) framework that unifies these approaches with linear attention (Section~\ref{subsec:linear-attn}).

\subsubsection{Structured State Space Models}
\label{subsubsec:s4}

\paragraph{S4: Structured State Spaces for Sequence Modeling.}
Gu et al.~\cite{gu2022efficiently} introduced S4, which discretizes the continuous-time linear dynamical system:
\begin{equation}
\dot{x}(t) = A\, x(t) + B\, u(t), \quad y(t) = C\, x(t),
\label{eq:ssm-continuous}
\end{equation}
where $A \in \mathbb{R}^{d_s \times d_s}$ is the state transition matrix, $B \in \mathbb{R}^{d_s \times 1}$ is the input matrix, $C \in \mathbb{R}^{1 \times d_s}$ is the output matrix, and $x(t) \in \mathbb{R}^{d_s}$ is the hidden state. Two innovations make this system practical for sequence modeling. First, the \emph{HiPPO} (High-order Polynomial Projection Operators) initialization sets $A$ to a specific structured matrix that ensures the hidden state $x(t)$ optimally approximates the input history $u(\tau)$ for $\tau \leq t$ in a Legendre polynomial basis. This initialization provides a theoretical guarantee of long-range memory that is absent from random initialization. Second, the \emph{DPLR} (Diagonal Plus Low-Rank) parameterization decomposes $A$ as $A = \Lambda + PQ^\top$ where $\Lambda$ is diagonal, enabling $O(N \log N)$ training via fast convolution in the frequency domain.

After discretization with step size $\Delta$ (using the zero-order hold method), the continuous system becomes:
\begin{equation}
x_t = \bar{A}\, x_{t-1} + \bar{B}\, u_t, \quad y_t = C\, x_t,
\label{eq:ssm-discrete}
\end{equation}
where $\bar{A} = \exp(\Delta A)$ and $\bar{B} = (\exp(\Delta A) - I) A^{-1} B$. This discrete recurrence supports three computation modes: (1) \emph{recurrent} mode ($O(d_s)$ per step) for inference, (2) \emph{convolutional} mode ($O(N \log N)$) for training via FFT, and (3) \emph{parallel scan} mode ($O(N d_s)$) for GPU-efficient training. S4 achieved breakthrough results on the Long Range Arena benchmark, demonstrating for the first time that a linear dynamical system with principled initialization can capture dependencies over sequences of thousands of tokens.

A critical limitation of S4 is that its parameters $A$, $B$, $C$ are \emph{time-invariant}---they do not depend on the input $u_t$. This means that S4 processes all tokens identically, regardless of their content, and cannot perform content-based selection or filtering. This limitation motivated the development of selective SSMs.

\subsubsection{Selective State Space Models: The Mamba Family}
\label{subsubsec:mamba}

\paragraph{Mamba.}
Gu and Dao~\cite{gu2024mamba} proposed Mamba, which makes the SSM parameters \emph{input-dependent} (selective):
\begin{equation}
x_t = \bar{A}_t\, x_{t-1} + \bar{B}_t\, u_t, \quad y_t = C_t\, x_t,
\label{eq:mamba}
\end{equation}
where $B_t = \text{Linear}_B(u_t)$, $C_t = \text{Linear}_C(u_t)$, and $\Delta_t = \text{softplus}(\text{Linear}_\Delta(u_t))$ are all functions of the input. The discretized parameters $\bar{A}_t = \exp(\Delta_t A)$ and $\bar{B}_t$ inherit the input dependence through $\Delta_t$. This selectivity mechanism allows the model to dynamically decide what to remember and what to forget at each time step: a large $\Delta_t$ causes $\bar{A}_t \approx 0$ (forgetting the state) and $\bar{B}_t \approx I$ (writing the input), while a small $\Delta_t$ causes $\bar{A}_t \approx I$ (preserving the state) and $\bar{B}_t \approx 0$ (ignoring the input).

The input dependence of the parameters breaks the time-invariance that enabled S4's convolutional computation mode, requiring a different training algorithm. Mamba introduces a \emph{hardware-aware parallel scan} that exploits the GPU memory hierarchy: the state $x_t$ is maintained in fast SRAM during the scan, avoiding the costly HBM reads/writes that would otherwise dominate the computation. The Mamba block combines the selective SSM with a SiLU-gated linear layer (analogous to the SwiGLU FFN in Transformers), creating a unified token-mixing and channel-mixing module. At 3B parameters, Mamba matches or exceeds Transformer quality on language modeling while achieving 5$\times$ higher inference throughput for long sequences.

\paragraph{Mamba-2 and the Structured State Space Duality (SSD).}
Dao and Gu~\cite{dao2024transformers} established a fundamental theoretical connection between SSMs and attention. When the state transition matrix is restricted to a scalar times the identity ($A = a \cdot I$), the SSM recurrence becomes equivalent to a \emph{semi-separable matrix} multiplication, which in turn generalizes causal linear attention:
\begin{equation}
y_i = \sum_{j \leq i} \left(\prod_{l=j+1}^{i} a_l\right) (C_i^\top B_j)\, u_j = \sum_{j \leq i} M_{ij}\, u_j,
\label{eq:ssd}
\end{equation}
where $M_{ij} = (\prod_{l=j+1}^{i} a_l)(C_i^\top B_j)$ is a semi-separable matrix. This \emph{Structured State Space Duality} (SSD) proves that SSMs, linear attention, and semi-separable matrices are three mathematically equivalent views of the same computation. The practical consequence is that Mamba-2 can be implemented using matrix multiplications (which map efficiently to GPU tensor cores) rather than sequential scans, achieving 2--8$\times$ speedup over Mamba-1 while maintaining identical expressivity. The SSD framework also introduces a multi-head SSM design analogous to multi-head attention, where each head maintains an independent state with shared $A$ parameters.

\paragraph{Mamba-3.}
Gu and Dao~\cite{gu2025mamba3} addressed two limitations of Mamba-2. First, the restriction to real-valued diagonal $A$ matrices means that Mamba-2 can only model exponential decay, not oscillatory dynamics. Mamba-3 restores \emph{complex-valued} SSM parameters (as in the original S4) through a conjugate-pair parameterization that enables efficient real-arithmetic implementation: each complex eigenvalue $a + bi$ is stored as a pair of real numbers, and the complex state update is computed using only real multiplications. Second, Mamba-3 introduces \emph{MIMO} (Multi-Input Multi-Output) SSMs where $B \in \mathbb{R}^{d_s \times D}$ and $C \in \mathbb{R}^{D \times d_s}$ (with $D$ being the number of channels), enabling cross-channel interaction within the SSM layer---a capability absent from the SISO (Single-Input Single-Output) formulation of Mamba-1/2. Third, Mamba-3 replaces the zero-order hold (ZOH) discretization with an \emph{exponential-trapezoidal} rule that provides a more accurate discretization for input-dependent parameters. These improvements establish a new Pareto frontier on the quality--efficiency trade-off, surpassing both Mamba-2 and Transformers of comparable size.

\paragraph{Gated Delta Networks.}
Yang et al.~\cite{yang2024gated} combined the delta rule (Section~\ref{subsubsec:delta-linear}) with the Mamba-2 SSD framework, yielding a state update that supports both global gating and key-specific memory operations:
\begin{equation}
S_t = \alpha_t \odot S_{t-1} - \beta_t\, k_t (k_t^\top S_{t-1}) + \beta_t\, k_t\, v_t^\top,
\label{eq:gated-delta-ssm}
\end{equation}
where $\alpha_t$ is a data-dependent decay gate and $\beta_t$ is the delta-rule learning rate. The first term provides global exponential decay (as in Mamba), the second term erases the old value associated with key $k_t$, and the third term writes the new value. This combination substantially improves associative recall compared to either mechanism alone, and the Householder WY representation enables efficient parallel training. Gated Delta Networks achieve the strongest results among sub-quadratic token mixers on in-context learning benchmarks, approaching the performance of full softmax attention.

\subsubsection{Modern RNN Variants}
\label{subsubsec:modern-rnn}

In parallel with the SSM line of research, several groups have revisited classical RNN architectures with modern design principles, demonstrating that well-designed recurrent models can be competitive with both Transformers and SSMs.

\paragraph{RWKV.}
Peng et al.~\cite{peng2023rwkv} introduced RWKV (Receptance Weighted Key Value), which implements channel-wise linear recurrence with learnable time-decay factors. The core computation is:
\begin{equation}
\text{wkv}_t = \frac{\sum_{j \leq t} e^{-(t-j)w + k_j} v_j}{\sum_{j \leq t} e^{-(t-j)w + k_j}},
\label{eq:rwkv}
\end{equation}
where $w > 0$ is a channel-wise decay factor and $k_j$ is the key at position $j$. This can be computed recursively in $O(d)$ per step, and the exponential weighting provides a soft recency bias. RWKV scales to 14B parameters with quality competitive with similarly-sized Transformers, demonstrating the viability of pure recurrence at production scale. RWKV-5/6 (Eagle and Finch)~\cite{peng2024eagle} extend the state from vector-valued to \emph{matrix-valued} representations ($S_t \in \mathbb{R}^{d_k \times d_v}$), increasing the state capacity from $O(d)$ to $O(d_k d_v)$ and enabling richer associative storage. RWKV-7~\cite{peng2025rwkv7} introduces the \emph{generalized delta rule} for state updates---the same ``erase-then-write'' mechanism as DeltaNet (Eq.~\ref{eq:deltanet})---combined with vector-valued (per-dimension) gating and a data-dependent in-context learning rate $\alpha_t$. This convergence of RWKV-7 and DeltaNet/Gated Delta Networks on the same delta-rule update mechanism, arrived at independently from different research traditions, provides strong evidence for the optimality of this state update rule.

\paragraph{RetNet.}
Sun et al.~\cite{sun2023retentive} proposed the \emph{retention} mechanism, which combines exponential decay with multi-scale heads:
\begin{equation}
\text{Retention}(Q, K, V)_n = \sum_{m \leq n} \gamma^{n-m} (Q_n^\top K_m)\, V_m,
\label{eq:retnet}
\end{equation}
where $\gamma \in (0, 1)$ is a per-head decay factor. RetNet supports three computation modes: parallel (matrix form for training), recurrent ($O(d^2)$ per step for inference), and chunk-wise (hybrid for long-sequence training). The multi-scale design assigns different decay rates to different heads, enabling simultaneous modeling of short-range and long-range dependencies. RetNet can be viewed as a special case of GLA (Eq.~\ref{eq:gla}) where the gate is a fixed scalar rather than a data-dependent matrix.

\paragraph{xLSTM.}
Beck et al.~\cite{beck2024xlstm} modernized the classical LSTM with two key innovations. \emph{sLSTM} (scalar LSTM) introduces \emph{exponential gating}: the input and forget gates use $\exp(\cdot)$ rather than $\sigma(\cdot)$, providing a much larger dynamic range that enables more precise memory control. Numerical stability is maintained through log-space computation. \emph{mLSTM} (matrix LSTM) replaces the scalar cell state with a \emph{matrix-valued memory} $C_t \in \mathbb{R}^{d \times d}$:
\begin{equation}
C_t = f_t \cdot C_{t-1} + i_t \cdot v_t\, k_t^\top, \quad h_t = \frac{C_t\, q_t}{\max(|n_t^\top q_t|, 1)},
\label{eq:mlstm}
\end{equation}
where $f_t = \exp(\tilde{f}_t)$ and $i_t = \exp(\tilde{i}_t)$ are exponential gates, and $n_t = f_t n_{t-1} + i_t k_t$ is a normalizer. The matrix memory update $C_t = f_t C_{t-1} + i_t v_t k_t^\top$ is structurally identical to gated linear attention (Eq.~\ref{eq:gla}), with the exponential gates providing a richer parameterization than sigmoid gates. mLSTM is fully parallelizable via the same chunk-wise algorithms used for GLA, demonstrating that classical recurrent architectures remain competitive when equipped with modern design principles and training algorithms.

\subsubsection{Long Convolution Architectures}
\label{subsubsec:long-conv}

A complementary approach to sub-quadratic sequence modeling replaces attention with long convolutions, which can be computed in $O(N \log N)$ via FFT.

\paragraph{Hyena.}
Poli et al.~\cite{poli2023hyena} proposed replacing attention with a hierarchy of long convolutions parameterized by implicit neural representations. Each Hyena layer applies a sequence of element-wise gating operations interleaved with long convolutions:
\begin{equation}
y = x_1 \odot (h_1 * (x_2 \odot (h_2 * (\cdots)))),
\label{eq:hyena}
\end{equation}
where $h_i$ are implicitly parameterized convolution filters (generated by a small MLP that maps position indices to filter values) and $x_i$ are data-dependent projections of the input. The implicit parameterization allows the filter length to be decoupled from the number of parameters, enabling arbitrarily long convolutions with a fixed parameter budget. Hyena achieves sub-quadratic complexity ($O(N \log N)$ via FFT) while matching attention quality on sequences up to 8K tokens.

\paragraph{MEGA.}
Ma et al.~\cite{ma2023mega} combined an \emph{exponential moving average} (EMA) with gated attention in a single module. The EMA component provides efficient local smoothing with learnable multi-scale decay rates, while a single-head attention mechanism (with $O(N^2)$ cost but small constant factor due to the single head) captures global dependencies. The combination achieves strong results on both language modeling and long-range benchmarks, demonstrating that even a minimal attention component can substantially improve the quality of recurrence-based models.

\subsubsection{Comparison and Analysis}

Table~\ref{tab:ssm-comparison} provides a comprehensive comparison of SSM and RNN-like token mixers along the dimensions of state update mechanism, state capacity, computational complexity, and selectivity.

\begin{table}[t]
\centering
\caption{Comparison of SSM and RNN-like token mixers. $N$: sequence length, $d_s$: state dimension, $d$: model dimension. ``Selective'' indicates whether the state update depends on the input content.}
\label{tab:ssm-comparison}
\small
\begin{tabular}{lcccccc}
\toprule
\textbf{Method} & \textbf{Year} & \textbf{State Update} & \textbf{State Size} & \textbf{Train} & \textbf{Infer/step} & \textbf{Selective} \\
\midrule
S4             & 2021 & Linear (HiPPO) & $d_s$ & $O(N \log N)$ & $O(d_s)$ & No \\
Mamba          & 2024 & Selective SSM & $d_s$ & $O(Nd_s)$ & $O(d_s)$ & Yes \\
Mamba-2 (SSD)  & 2024 & Scalar-$A$ SSM & $d_s$ & $O(Nd_s)$ & $O(d_s)$ & Yes \\
Mamba-3        & 2025 & Complex MIMO SSM & $d_s$ & $O(Nd_s)$ & $O(d_s)$ & Yes \\
Gated DeltaNet & 2024 & Gate + delta rule & $d_k {\times} d_v$ & $O(Nd_k d_v)$ & $O(d_k d_v)$ & Yes \\
RWKV           & 2023 & Exp-decay recurrence & $d$ & $O(Nd)$ & $O(d)$ & Partial \\
RWKV-5/6       & 2024 & Matrix-valued state & $d_k {\times} d_v$ & $O(Nd_k d_v)$ & $O(d_k d_v)$ & Partial \\
RWKV-7         & 2025 & Generalized delta rule & $d_k {\times} d_v$ & $O(Nd_k d_v)$ & $O(d_k d_v)$ & Yes \\
RetNet         & 2023 & Exp-decay retention & $d_k {\times} d_v$ & $O(Nd_k d_v)$ & $O(d_k d_v)$ & No \\
xLSTM (mLSTM)  & 2024 & Exp-gated matrix memory & $d {\times} d$ & $O(Nd^2)$ & $O(d^2)$ & Yes \\
Hyena          & 2023 & Long convolution & --- & $O(N \log N)$ & $O(N)$ & Partial \\
MEGA           & 2023 & EMA + single-head attn & $d$ & $O(N^2/H{+}Nd)$ & $O(N{+}d)$ & Partial \\
\bottomrule
\end{tabular}
\end{table}

The SSD framework~\cite{dao2024transformers} provides a unifying theoretical lens for understanding the relationships among these methods. Under SSD, the key distinction is the \emph{state update rule}: vanilla linear attention and RetNet use additive updates ($S_t = \alpha S_{t-1} + k_t v_t^\top$), GLA and mLSTM use gated updates ($S_t = G_t \odot S_{t-1} + k_t v_t^\top$), and DeltaNet/RWKV-7/Gated Delta Networks use delta-rule updates ($S_t = S_{t-1} - \beta_t k_t(k_t^\top S_{t-1}) + \beta_t k_t v_t^\top$). This progression---from additive to gated to delta-rule---represents an increasingly expressive hierarchy of state update mechanisms, each addressing a specific limitation of its predecessor: gating solves the monotonic accumulation problem, and the delta rule solves the associative interference problem.

The fundamental limitation shared by all SSM and RNN-like models is the \emph{state bottleneck}: the fixed-size state $S \in \mathbb{R}^{d_k \times d_v}$ can store at most $O(d_k d_v)$ bits of information, regardless of the sequence length. This bottleneck manifests as degraded performance on tasks requiring precise retrieval from long contexts---particularly in-context learning, needle-in-a-haystack, and multi-hop reasoning---where softmax attention's ability to access any token in the KV cache with $O(1)$ precision provides an insurmountable advantage. The dominant strategy for addressing this limitation is \emph{hybrid architectures} (Section~\ref{subsec:hybrid}) that interleave a small number of full softmax attention layers (for precise recall) with SSM or linear attention layers (for efficient long-range processing). The independent convergence of Jamba~\cite{lieber2024jamba_tr} and MiniMax-01~\cite{minimax2025} on a ${\sim}$7:1 SSM-to-attention ratio suggests a natural balance point, though the optimal ratio likely depends on the task distribution and model scale.

Future directions include more powerful selectivity mechanisms (e.g., complex-valued dynamics in Mamba-3, MIMO interactions), the development of state update rules that can match softmax attention's recall capability without the quadratic cost, and the co-design of SSM architectures with hardware accelerators to close the gap between theoretical and realized throughput. The convergence of RWKV-7, DeltaNet, and Gated Delta Networks on the delta-rule update mechanism, arrived at independently from different research traditions, provides strong evidence that this update rule represents a fundamental optimum for linear-complexity sequence modeling.
